\begin{abstract}
Retrieval-augmented generation (RAG) systems depend critically on how documents are chunked and searched. Fine-grained chunks can improve retrieval precision but expand the search space, increasing latency and cost; larger chunks reduce the number of candidates but make dense similarity less reliable, as the representation for each chunk mixes multiple topics and introduces more semantic noise. This trade-off becomes especially limiting in deep research tasks, where retrieval must be both fast and precise across large, heterogeneous corpora. We introduce \textbf{\name{}}, a metadata-guided retrieval framework that uses topic-level signals as a semantic compass for selecting relevant evidence. Instead of relying only on cosine similarity between queries and noisy chunk embeddings, \name{} enriches chunk representations with topic metadata in the same embedding space and trains a lightweight retriever through LLM-teacher distillation. At inference time, \name{} performs topic-aware retrieval without additional LLM calls, improving both efficiency and evidence quality. Across six complex retrieval benchmarks, \name{} improves information efficiency (IE) by \textbf{8.24\%} on average with over \textbf{5$\times$ lower latency} than the strongest efficient RAG baselines\footnote{
Code is available on \href{https://github.com/AmirAbaskohi/MCompassRAG}
{{\textcolor{black}{\faGithub}}~GitHub}.
}. 
\end{abstract}